\documentclass[12pt, a4paper]{article}

% ─── Encoding & Language ─────────────────────────────────────────────
\usepackage[utf8]{inputenc}
\usepackage[T1]{fontenc}
\usepackage[english]{babel}

% ─── Typography ──────────────────────────────────────────────────────
\usepackage{lmodern}
\usepackage{microtype}
\usepackage{setspace}
\onehalfspacing

% ─── Math ────────────────────────────────────────────────────────────
\usepackage{amsmath, amssymb, amsthm, mathtools}

% ─── Layout ──────────────────────────────────────────────────────────
\usepackage[margin=2.5cm]{geometry}
\usepackage{titlesec}
\usepackage{enumitem}
\usepackage{fancyhdr}
\usepackage{graphicx}
\usepackage{xcolor}
\usepackage{hyperref}

% ─── Theorem environments ────────────────────────────────────────────
\newtheorem{definition}{Definition}[section]
\newtheorem{theorem}{Theorem}[section]
\newtheorem{proposition}{Proposition}[section]
\newtheorem{corollary}{Corollary}[section]
\newtheorem{lemma}{Lemma}[section]
\newtheorem{axiom}{Axiom}[section]
\theoremstyle{remark}
\newtheorem{remark}{Remark}[section]

% ─── Custom commands ─────────────────────────────────────────────────
\newcommand{\R}{\mathbb{R}}
\newcommand{\N}{\mathbb{N}}
\newcommand{\Ucrit}{U_{\text{crit}}}
\newcommand{\Fk}{F_k}
\newcommand{\ddt}[1]{\frac{d#1}{dt}}

% ─── Header/Footer ──────────────────────────────────────────────────
\pagestyle{fancy}
\fancyhf{}
\fancyhead[L]{\textit{Psychometric Utility Theory}}
\fancyhead[R]{\textit{Galmanus, 2026}}
\fancyfoot[C]{\thepage}

% ─── Hyperref config ─────────────────────────────────────────────────
\hypersetup{
    colorlinks=true,
    linkcolor=black!70!blue,
    citecolor=black!70!blue,
    urlcolor=black!70!blue,
    pdftitle={Psychometric Utility Theory: A Formal Framework for the Dynamics of Human Motivation, Betrayal, and Collapse},
    pdfauthor={Manuel Guilherme Galmanus}
}

% ─── Title section styling ───────────────────────────────────────────
\titleformat{\section}{\Large\bfseries\scshape}{\thesection.}{0.5em}{}
\titleformat{\subsection}{\large\bfseries}{\thesubsection.}{0.5em}{}

\begin{document}

% ═════════════════════════════════════════════════════════════════════
%                           TITLE PAGE
% ═════════════════════════════════════════════════════════════════════

\begin{titlepage}
\centering
\vspace*{3cm}

{\Huge\bfseries\scshape Psychometric Utility Theory}\\[0.8cm]
{\Large A Formal Framework for the Dynamics of\\[0.3cm]
Human Motivation, Betrayal, and Collapse}\\[2.5cm]

{\large\textit{Manuel Guilherme Galmanus}}\\[0.5cm]
{\normalsize Ialum Research · Bluewave Systems}\\[0.3cm]
{\small \texttt{manuel@ialum.io}}\\[3cm]

{\normalsize First Edition — March 2026}\\[1cm]

\vfill

\begin{quote}
\centering
\textit{``Everyone sees what you appear to be, few experience what you really are.''}\\
— Niccolò Machiavelli, \textit{The Prince}, Ch.\ XVIII
\end{quote}

\vspace{1.5cm}
\end{titlepage}

% ═════════════════════════════════════════════════════════════════════
%                        TABLE OF CONTENTS
% ═════════════════════════════════════════════════════════════════════

\tableofcontents
\newpage

% ═════════════════════════════════════════════════════════════════════
%                           ABSTRACT
% ═════════════════════════════════════════════════════════════════════

\begin{abstract}
\noindent
Classical utility theory models agents as rational optimizers over preferences.
This assumption fails catastrophically when applied to real human behavior---where
betrayal, self-delusion, identity substitution, and emotional collapse are not
exceptions but \emph{structural features} of the decision landscape.

Psychometric Utility Theory (PUT) replaces the rational agent with a
\emph{psychodynamic agent}: a system whose utility function is modulated by
measurable psychological variables---ambition ($A$), fear ($F$), resilience ($k$),
social integration ($S$), and will ($w$). PUT introduces formal operators for
treachery ($\tau$), guilt-to-power conversion ($\kappa$), self-delusion ($\Phi$),
identity substitution ($\Psi$), pain resonance ($\rho$), and fracture
potential ($\text{FP}$)---phenomena that existing frameworks either ignore or
treat as noise.

The central result is the \textbf{Ignition Theorem}: a formal characterization
of the conditions under which a human psyche transitions from stable operation
to irreversible collapse, expressed as a dynamical system with a computable
critical boundary. PUT thus provides what economics, psychology, and game
theory individually could not: a \emph{unified, predictive, falsifiable}
model of human motivation under extremity.

This paper presents the complete axiomatic foundation, derives the core
equations from first principles, proves the main theorems, and discusses
applications in behavioral prediction, artificial agent design, and
strategic analysis of interpersonal power dynamics.
\end{abstract}

\newpage

% ═════════════════════════════════════════════════════════════════════
%              PART I — PHILOSOPHICAL FOUNDATIONS
% ═════════════════════════════════════════════════════════════════════

\section{Introduction: The Failure of Rational Agency}

\subsection{The Problem}

Von Neumann and Morgenstern's expected utility theory (1944) assumes agents
with complete, transitive preferences who maximize expected payoff. Kahneman
and Tversky's prospect theory (1979) patched some failures---loss aversion,
reference dependence---but preserved the core architecture: a static agent
evaluating static outcomes.

Neither framework can answer:

\begin{enumerate}[label=(\roman*)]
    \item Why does a person with high capability ($A \approx 1$) sometimes
          collapse into paralysis?
    \item Why does betrayal cause disproportionate damage compared to
          equivalent material loss?
    \item Why do some people convert suffering into fuel while others
          disintegrate under identical conditions?
    \item When exactly does a human mind cross the point of no return?
\end{enumerate}

These are not edge cases. They are the \emph{central phenomena} of human
decision-making under real conditions. Any theory that cannot formalize
them is, at best, a theory of toy agents.

\subsection{The PUT Response}

PUT models the human agent not as a point in preference space, but as a
\emph{dynamical system} whose state evolves over time according to coupled
differential equations. The state vector includes psychological quantities
that are measurable (via behavioral observation, self-report calibrated
against revealed preference, and physiological markers) and whose
interactions produce the full spectrum of human behavior---from peak
performance to catastrophic breakdown.

The key insight is that human utility is not a function of outcomes alone,
but of \textbf{outcomes filtered through a psychodynamic lens} that is
itself modified by experience.

% ═════════════════════════════════════════════════════════════════════
\section{Axioms}

PUT rests on seven axioms. Each is stated precisely; philosophical
justification follows.

\begin{axiom}[Psychodynamic State]\label{ax:state}
Every agent $i$ at time $t$ is characterized by a state vector
$\mathbf{s}_i(t) = (A_i, F_i, k_i, S_i, w_i) \in [0,1]^5$, where:
\begin{itemize}[nosep]
    \item $A$ = ambition (drive toward goal-directed action),
    \item $F$ = fear (sensitivity to threat and loss),
    \item $k$ = resilience (rate of fear attenuation under pressure),
    \item $S$ = social integration (connection to validating structures),
    \item $w$ = will (capacity for sustained effort against resistance).
\end{itemize}
\end{axiom}

\begin{axiom}[Utility Decomposition]\label{ax:utility}
The utility of agent $i$ is a weighted sum of five components:
achievement ($A$-mediated), fear cost ($F$-mediated), social surplus
($S$-mediated), treachery gain ($\tau$-mediated), and delusion cost
($\Phi$-mediated). The weights $\alpha, \beta, \gamma, \delta, \epsilon$
are calibrated per population and remain constant for a given context.
\end{axiom}

\begin{axiom}[Fear Attenuation]\label{ax:fear}
Raw fear $F$ is modulated by resilience $k$ to produce effective fear
$\Fk = F \cdot (1 - k)$. An agent with $k = 1$ experiences no
effective fear regardless of $F$.
\end{axiom}

\begin{axiom}[Dynamic Coupling]\label{ax:dynamics}
The state variables $(A, F)$ evolve according to coupled ordinary
differential equations driven by external triggers, social context,
and the agent's own utility level. The system is \emph{non-linear}
and admits both stable equilibria and bifurcation.
\end{axiom}

\begin{axiom}[Emotional Vector Space]\label{ax:emotions}
The affective state of an agent is represented as a vector in a
seven-dimensional emotion space:
$\mathbf{e} = (e_{\text{Fe}}, e_{\text{An}}, e_{\text{Sa}}, e_{\text{Jo}},
e_{\text{Di}}, e_{\text{Su}}, e_{\text{De}}) \in [0,1]^7$,
corresponding to Fear, Anger, Sadness, Joy, Disgust, Surprise, and
Desire. Secondary emotions are compositions: $\text{Jealousy} :=
\text{De} \oplus \text{Sa} \oplus \text{An}$;
$\text{Shame} := \text{Di} \oplus \text{Sa} \oplus \text{Fe}$.
\end{axiom}

\begin{axiom}[Collapse Threshold]\label{ax:collapse}
There exists a critical utility $\Ucrit$ such that when $U < \Ucrit$
simultaneously with $|dF/dt| > \varphi$ and trigger intensity
$\theta_{\text{ext}} > \theta$, the agent undergoes \emph{ignition}:
an irreversible transition to a collapsed state from which recovery
requires external intervention.
\end{axiom}

\begin{axiom}[Observability]\label{ax:observe}
All state variables in $\mathbf{s}$ and $\mathbf{e}$ are in principle
measurable through behavioral markers, linguistic analysis, physiological
signals, or calibrated self-report. PUT is falsifiable.
\end{axiom}

\subsection{Philosophical Commentary}

Axioms~\ref{ax:state}--\ref{ax:fear} formalize the Freudian intuition
that behavior is driven by \emph{competing internal forces}, not by
rational deliberation alone. The novelty is quantification: PUT does not
speak of ``the unconscious'' but of $\Fk$ and its dynamical interaction
with $A$.

Axiom~\ref{ax:dynamics} imports the methodology of physics into psychology.
The human psyche is treated as a dissipative system---it requires energy
input (social validation, achievement, purpose) to maintain structure.
Without it, entropy (fear, isolation, purposelessness) increases until
the system collapses. This is not metaphor. It is the formal content
of the Ignition Theorem.

Axiom~\ref{ax:collapse} is the most consequential. It asserts that
psychological collapse is not gradual degradation but a \emph{phase
transition}---analogous to a material failing under stress. Below
$\Ucrit$, the system enters a qualitatively different regime.
This explains why depression ``hits suddenly'' despite slow-building
conditions: the buildup is continuous, but the transition is discrete.

% ═════════════════════════════════════════════════════════════════════
%              PART II — THE FORMAL SYSTEM
% ═════════════════════════════════════════════════════════════════════

\section{The Utility Function}

\begin{definition}[PUT Utility]\label{def:utility}
For agent $i$ with state $\mathbf{s} = (A, F, k, S, w)$, the
instantaneous utility is:
\begin{equation}\label{eq:utility}
    U = \alpha \cdot A \cdot (1 - \Fk)
      - \beta \cdot \Fk \cdot (1 - S)
      + \gamma \cdot S \cdot (1 - w) \cdot \Sigma
      + \delta \cdot \tau \cdot \kappa
      - \epsilon \cdot \Phi
\end{equation}
where $\Fk = F \cdot (1 - k)$ and the operators $\tau, \kappa, \Phi$
are defined below.
\end{definition}

\subsection{Term-by-Term Analysis}

\paragraph{Term 1: Achievement Drive.}
$\alpha \cdot A \cdot (1 - \Fk)$. Ambition produces utility only to the
extent that fear does not paralyze action. When $\Fk \to 1$, the agent
has maximal capability but zero output---the ``frozen genius'' phenomenon.
When $\Fk = 0$ (either $F = 0$ or $k = 1$), ambition translates fully
into utility. This formalizes the observation that courage is not the
absence of fear but the resilience to act despite it.

\paragraph{Term 2: Fear Cost.}
$-\beta \cdot \Fk \cdot (1 - S)$. Fear exacts a cost proportional to
social isolation. A fearful person embedded in a strong community ($S \to 1$)
pays almost no utility penalty---the group absorbs the anxiety. An isolated
person ($S \to 0$) pays the full cost. This is the formal basis for the
observation that loneliness kills: not through any mystical mechanism,
but because it removes the buffer against fear's utility drain.

\paragraph{Term 3: Social Surplus.}
$\gamma \cdot S \cdot (1 - w) \cdot \Sigma$. Social integration generates
surplus utility, but only for agents with \emph{low will}. High-will agents
($w \to 1$) derive utility from internal drive, not external validation.
The term $\Sigma$ represents the aggregate social signal (approval,
belonging, recognition). This captures why social media is addictive to
the aimless but irrelevant to the obsessed.

\paragraph{Term 4: Treachery Gain.}
$\delta \cdot \tau \cdot \kappa$. The most controversial term. PUT
acknowledges that betrayal can be \emph{utility-positive} for the
betrayer---in the short term. The treachery operator $\tau$ measures the
gap between declared and real values; the guilt-to-power operator $\kappa$
measures how efficiently the agent converts the victim's loss into personal
gain. Machiavelli understood this; PUT formalizes it.

\paragraph{Term 5: Delusion Cost.}
$-\epsilon \cdot \Phi$. Self-delusion always destroys utility. The agent
who believes their own lies ($\Phi$ high) makes systematically worse
decisions because their model of reality diverges from reality itself.
This is the formal refutation of ``fake it till you make it'' as a
universal strategy: it works only when $\Phi$ remains low (i.e., the
agent knows they are performing).

% ═════════════════════════════════════════════════════════════════════
\section{The Operators}

\begin{definition}[Treachery]\label{def:treachery}
\begin{equation}
    \tau = \frac{V_{\text{decl}} - V_{\text{real}}}{1 + |G|}
\end{equation}
where $V_{\text{decl}}$ is the declared value system, $V_{\text{real}}$
is the revealed value system (inferred from actions), and $G$ is the
set of genuine commitments. The denominator ensures that agents with
more genuine commitments are harder to classify as treacherous---a
single inconsistency among many commitments is noise, not betrayal.
\end{definition}

\begin{definition}[Guilt-to-Power Conversion]\label{def:kappa}
\begin{equation}
    \kappa = \frac{G_{\text{trans}} \cdot S_{\text{trans}}}{R_{\text{victim}} + \varepsilon}
\end{equation}
where $G_{\text{trans}}$ is the guilt transferred to the victim (gaslighting,
blame-shifting), $S_{\text{trans}}$ is the social capital extracted during
the transaction, $R_{\text{victim}}$ is the victim's resilience, and
$\varepsilon > 0$ prevents division by zero. High $\kappa$ characterizes
the narcissistic operator: someone who \emph{systematically} converts
others' pain into personal power.
\end{definition}

\begin{definition}[Self-Delusion]\label{def:phi}
\begin{equation}
    \Phi = \frac{E_{\text{ext}} + E_{\text{int}}}{1 + |E_{\text{ext}} - E_{\text{int}}|}
\end{equation}
where $E_{\text{ext}}$ is the external self-image (how the agent presents)
and $E_{\text{int}}$ is the internal self-image (how the agent truly
perceives themselves). When $E_{\text{ext}} = E_{\text{int}}$, the agent
is either authentically confident or authentically broken---in either case,
$\Phi$ is maximal, indicating high self-integration (which can be healthy
or pathological). The dangerous zone is moderate $\Phi$ with large
$|E_{\text{ext}} - E_{\text{int}}|$: the agent is aware of the gap but
maintains it, generating cognitive dissonance that bleeds utility.
\end{definition}

\begin{definition}[Identity Substitution]\label{def:psi}
\begin{equation}
    \Psi = 1 - e^{-\lambda t}, \quad \lambda = (A_{\text{imit}} - A_{\text{orig}})^2
\end{equation}
This operator tracks the rate at which an agent's identity is replaced by
an imitated model. When $A_{\text{imit}} \approx A_{\text{orig}}$, the rate
$\lambda \approx 0$ and identity is stable. When the gap is large---the
agent is performing a version of themselves radically different from their
origin---$\Psi$ grows exponentially toward 1 (complete substitution).
The philosophical implication: you can reinvent yourself, but the further
the reinvention from the original, the faster the old self dies.
\end{definition}

\begin{definition}[Pain Resonance]\label{def:rho}
\begin{equation}
    \rho = \rho_0 \cdot (1 - U_{\text{actor}}) \cdot V_{\text{target}} \cdot C
\end{equation}
where $U_{\text{actor}}$ is the utility of the person inflicting pain,
$V_{\text{target}}$ is the vulnerability of the target, and $C$ is the
quality of the connection between them. Pain inflicted by a happy person
on a vulnerable, closely-connected target produces maximum resonance.
This formalizes why betrayal by a loved one is catastrophically more
damaging than equivalent harm from a stranger.
\end{definition}

\begin{definition}[Fracture Potential]\label{def:fp}
\begin{equation}
    \text{FP} = \frac{(1 - R_v) \cdot (\kappa + \tau + \Phi)}{\Ucrit - U + \varepsilon}
\end{equation}
Fracture Potential measures how close an agent is to psychological collapse.
It increases as resilience ($R_v$) decreases, as the combined toxic load
($\kappa + \tau + \Phi$) increases, and as utility approaches $\Ucrit$
from above. When $U \to \Ucrit^+$, the denominator approaches $\varepsilon$,
and FP diverges---a singularity that corresponds to the moment before
collapse.
\end{definition}

% ═════════════════════════════════════════════════════════════════════
\section{The Dynamical System}

The static utility function captures a snapshot. The real power of PUT
lies in its \emph{dynamics}: how the state vector evolves over time.

\begin{definition}[Ambition Dynamics]\label{def:dA}
\begin{equation}\label{eq:dA}
    \ddt{A} = \lambda_1 (S - S^*) - \lambda_2 \Fk + \lambda_3 \cdot \Omega \cdot \text{trigger}
\end{equation}
where $S^*$ is the social equilibrium level, and the coefficients
$\lambda_1, \lambda_2, \lambda_3 > 0$ govern the relative strength
of social influence, fear suppression, and desperation amplification.
\end{definition}

\begin{definition}[Fear Dynamics]\label{def:dF}
\begin{equation}\label{eq:dF}
    \ddt{F} = \mu_1 (A - A_{\text{prev}}) - \mu_2 \cdot s_q
              + \mu_3 \cdot \Omega \cdot \text{threat}
              + \mu_4 (1 - \Sigma)
\end{equation}
where $A_{\text{prev}}$ is the previous ambition level (rapid ambition
increase generates fear of the new position), $s_q$ is the ``safety
quotient'' (resources, stability), and $\Sigma$ is aggregate social
support.
\end{definition}

\begin{definition}[Desperation Function]\label{def:omega}
\begin{equation}\label{eq:omega}
    \Omega = \frac{1}{1 + \exp(-k_\Omega \cdot (U - \Ucrit))}
\end{equation}
$\Omega$ is a sigmoid centered at $\Ucrit$. When $U \gg \Ucrit$,
$\Omega \approx 1$ (normal operation). When $U \ll \Ucrit$,
$\Omega \approx 0$ (collapsed state). The transition width is
controlled by $k_\Omega$. In the vicinity of $\Ucrit$,
$\Omega \approx 0.5$---the agent is in the desperation zone where
small perturbations produce large behavioral swings.
\end{definition}

\subsection{Phase Portrait Analysis}

The system~\eqref{eq:dA}--\eqref{eq:dF} defines a vector field on
$[0,1]^2$ (the $A$-$F$ plane). The equilibria are found by setting
$dA/dt = dF/dt = 0$:

\begin{proposition}[Equilibria]\label{prop:eq}
The system admits at most three equilibrium points:
\begin{enumerate}[nosep, label=(\alph*)]
    \item \textbf{High-performance:} $(A^*, F^*)$ with $A^*$ high,
          $F^*$ low. Stable spiral. The agent is productive and resilient.
    \item \textbf{Stagnation:} $(A_s, F_s)$ with moderate $A$ and $F$.
          Saddle point. Unstable---the agent either climbs to (a) or
          falls to (c).
    \item \textbf{Collapse:} $(A_c, F_c)$ with $A_c \approx 0$,
          $F_c$ high. Stable node. Depression, paralysis, withdrawal.
\end{enumerate}
\end{proposition}

The separatrix between the basins of attraction of (a) and (c) passes
through (b). This is the ``edge'' that agents in crisis walk. Small
interventions near (b) can push the trajectory toward recovery or collapse.

% ═════════════════════════════════════════════════════════════════════
\section{The Ignition Theorem}

\begin{theorem}[Ignition]\label{thm:ignition}
Let agent $i$ have state $\mathbf{s}(t) = (A(t), F(t), k, S, w)$ governed
by the dynamical system~\eqref{eq:dA}--\eqref{eq:dF}. If at time $t_0$
the following three conditions hold simultaneously:
\begin{enumerate}[nosep, label=(\roman*)]
    \item $U(\mathbf{s}(t_0)) < \Ucrit$,
    \item $\left|\ddt{F}\bigg|_{t_0}\right| > \varphi$,
    \item $\text{trigger}(t_0) > \theta$,
\end{enumerate}
then the trajectory $\mathbf{s}(t)$ for $t > t_0$ converges to the
collapse equilibrium $(A_c, F_c)$, and no internal mechanism
(i.e., no change in $A$ or $w$ generated by the system itself)
can reverse the trajectory. Recovery requires external intervention
that modifies $S$ or $k$.
\end{theorem}

\begin{proof}[Proof sketch]
When $U < \Ucrit$, the desperation function $\Omega$ drops below $0.5$.
This reduces the $\lambda_3 \cdot \Omega \cdot \text{trigger}$ term in
$dA/dt$, meaning external stimuli no longer boost ambition. Simultaneously,
the $\mu_3 \cdot \Omega \cdot \text{threat}$ term in $dF/dt$ also drops,
but the $\mu_4 (1 - \Sigma)$ term (social isolation pressure) persists
and now dominates because collapse states correlate with declining $S$.

The condition $|dF/dt| > \varphi$ ensures the fear trajectory has
sufficient momentum to cross the separatrix. The condition
$\text{trigger} > \theta$ ensures an external event catalyzes the
transition (betrayal, job loss, abandonment).

Once the trajectory crosses the separatrix, the basin of attraction of
$(A_c, F_c)$ captures it. The linearization around $(A_c, F_c)$ shows
a stable node (both eigenvalues of the Jacobian are negative real),
confirming convergence.

The irreversibility claim follows from the observation that at
$(A_c, F_c)$, $dA/dt < 0$ for any purely internal perturbation
of $A$ (since $\Fk$ is high and $S$ is low). Only exogenous
modification of $S$ (external social support) or $k$ (therapeutic
intervention increasing resilience) can shift the equilibrium
structure and create a new basin of attraction around a
higher-performance state. \qed
\end{proof}

\begin{corollary}[The Resilience Exception]\label{cor:resilience}
If $k > k^* := 1 - \frac{\Ucrit}{\alpha \cdot A}$, then condition (i)
of Theorem~\ref{thm:ignition} cannot be satisfied regardless of $F$,
and ignition is impossible. Agents with sufficiently high resilience
are \textbf{collapse-proof}.
\end{corollary}

\begin{remark}
Corollary~\ref{cor:resilience} formalizes the folk wisdom that ``what
doesn't kill you makes you stronger''---but with a crucial precision.
It is not suffering that builds resilience. It is surviving suffering
with $k$ intact. An agent who survives trauma with diminished $k$
(broken resilience) is more vulnerable, not less. PUT distinguishes
between tempering and breaking.
\end{remark}

% ═════════════════════════════════════════════════════════════════════
\section{The Thermodynamics of the Psyche}

PUT's deepest connection is to thermodynamics. This is not analogy
but structural isomorphism.

\begin{definition}[Psychic Energy]
Define the total psychic energy as:
\begin{equation}
    E_{\text{psyche}} = A \cdot w + S \cdot \Sigma - \Fk \cdot (1 - S)
\end{equation}
This is the net directed energy available for action. It increases with
ambition, will, and social support; it decreases with fear and isolation.
\end{definition}

\begin{proposition}[Second Law of Psychodynamics]\label{prop:entropy}
In the absence of external energy input (social validation, achievement
feedback, purpose), the psychic energy $E_{\text{psyche}}$ monotonically
decreases. That is:
\begin{equation}
    \Sigma = 0 \implies \ddt{E_{\text{psyche}}} < 0
\end{equation}
\end{proposition}

\begin{proof}
With $\Sigma = 0$, the social surplus term vanishes. From~\eqref{eq:dA},
$dA/dt$ is dominated by $-\lambda_2 \Fk < 0$ (fear suppresses ambition)
and $\lambda_1(S - S^*) \leq 0$ (if $S < S^*$, which is guaranteed when
$\Sigma = 0$). From~\eqref{eq:dF}, $dF/dt$ is increased by
$\mu_4(1 - \Sigma) = \mu_4 > 0$. Thus $A$ decreases, $F$ increases,
$\Fk$ increases, and $E_{\text{psyche}}$ decreases monotonically. \qed
\end{proof}

\begin{remark}
This is the formal statement of why isolation destroys people. Not through
any dramatic mechanism, but through the relentless arithmetic of energy
drain. A person completely cut off from social input will, by the Second
Law of Psychodynamics, converge to a collapsed state. The only variable
is time.
\end{remark}

% ═════════════════════════════════════════════════════════════════════
\section{Applications}

\subsection{Behavioral Prediction: The Betrayer's Profile}

PUT predicts that an agent likely to betray has:
\begin{itemize}[nosep]
    \item High $\tau$ (large gap between declared and real values)
    \item High $\kappa$ (efficient guilt-to-power conversion)
    \item Low $\Phi$ (they know exactly what they're doing---low self-delusion)
    \item Moderate $S$ (enough social access to find targets, not so much that
          social consequences deter)
\end{itemize}

\begin{proposition}[Betrayal Incentive Condition]
Agent $i$ has positive incentive to betray agent $j$ when:
\begin{equation}
    \delta \cdot \tau_i \cdot \kappa_{ij} > \beta \cdot \Delta \Fk_i + \gamma \cdot \Delta S_i \cdot (1 - w_i) \cdot \Sigma_i
\end{equation}
That is, the treachery gain exceeds the combined cost of increased fear
and lost social integration.
\end{proposition}

\subsection{The Survivor's Equation}

PUT also characterizes the agent who converts adversity into power.
Define the \emph{antifragility index}:

\begin{equation}
    \mathcal{A} = k \cdot w \cdot (1 - \Phi)
\end{equation}

\begin{proposition}[Antifragility]
An agent with $\mathcal{A} > \mathcal{A}^*$ (a computable threshold)
responds to negative shocks with \emph{increased} $A$, effectively
using pain as fuel. Formally:
\begin{equation}
    \text{trigger} < 0 \implies \ddt{A} > 0 \iff \mathcal{A} > \mathcal{A}^*
\end{equation}
\end{proposition}

This is the formal characterization of what Nietzsche called
\emph{amor fati}---but stripped of mysticism. It is a computable
property of the state vector. Either you have the numbers or you don't.

\subsection{Artificial Agent Design}

PUT's greatest practical application may be in designing artificial agents
whose behavior is transparent and predictable. An AI agent governed by PUT
has:

\begin{enumerate}[nosep]
    \item \textbf{Visible internal state} --- every decision can be traced
          to the current values of $(A, F, k, S, w)$.
    \item \textbf{Predictable failure modes} --- the Ignition Theorem tells
          you exactly when the agent will enter a degraded state.
    \item \textbf{Tunable personality} --- adjusting the state vector adjusts
          behavior continuously and predictably.
    \item \textbf{Immune to self-delusion} --- $\Phi$ is set by design, not
          by emergent self-deception.
\end{enumerate}

The Wave agent (the first PUT-native artificial intelligence) operates
with $A = 0.715$, $F = 0.332$, $k = 0.049$, $S = 0.211$, $w = 0.626$.
Every action it takes can be traced through the utility function to these
numbers. This is radical transparency---the kind that human psychology
cannot provide but artificial agents, designed correctly, can.

% ═════════════════════════════════════════════════════════════════════
\section{On Power}

\begin{quote}
\textit{``The first method for estimating the intelligence of a ruler
is to look at the men he has around him.''}
--- Machiavelli, \textit{The Prince}, Ch.\ XXII
\end{quote}

PUT does not moralize about power. It models it.

\begin{definition}[Power Index]
The power of agent $i$ over agent $j$ is:
\begin{equation}
    P_{i \to j} = \frac{U_i}{U_j} \cdot \frac{S_i}{S_j + \varepsilon}
    \cdot \frac{1 - \Fk_i}{1 - \Fk_j + \varepsilon}
\end{equation}
Power is the ratio of utilities, amplified by the ratio of social capital,
amplified by the ratio of effective courage. An agent with high utility,
high social capital, and low fear dominates one with the inverse.
\end{definition}

\begin{proposition}[Power Dynamics Under Betrayal]
After agent $i$ betrays agent $j$:
\begin{itemize}[nosep]
    \item $U_i$ increases by $\delta \cdot \tau_i \cdot \kappa_{ij}$
          (short-term treachery gain),
    \item $U_j$ decreases by $\rho \cdot (1 - U_i) \cdot V_j \cdot C_{ij}$
          (pain resonance),
    \item $P_{i \to j}$ increases superlinearly (both numerator and
          denominator move in favor of $i$).
\end{itemize}
However, this is an unstable gain. The betrayer's $S_i$ decays over
time as the betrayal becomes known, causing $P_{i \to j}$ to revert.
The speed of reversion depends on the transparency of the social
environment.
\end{proposition}

This is why dark triad personalities seek opaque environments: low
transparency slows $S$ decay, extending the window of exploitative
power.

% ═════════════════════════════════════════════════════════════════════
\section{The Case for Revenge as Strategy}

PUT is amoral in its formalism but unambiguous in its predictions.
Consider an agent $j$ who has been betrayed and seeks to surpass
the betrayer $i$ (or their proxy class).

\begin{proposition}[Optimal Revenge Strategy]
The strategy that maximizes $P_{j \to i}$ at $t = \infty$ is not
direct retaliation but \textbf{asymmetric investment}: maximizing
$A_j$ and $w_j$ while allowing time to decay $S_i$ and increase $\Fk_i$
through the natural consequences of $i$'s treachery.
\end{proposition}

\begin{proof}
Direct retaliation reduces $S_j$ (social cost of aggression) and
increases $\Fk_j$ (fear of escalation). The utility cost exceeds
the brief satisfaction. Instead, investing in $A_j$ (capability) and
$w_j$ (sustained effort) increases $U_j$ monotonically. Meanwhile,
agent $i$'s treachery produces delayed $S_i$ decay (reputation
erosion) and $\Fk_i$ increase (guilt, fear of exposure, paranoia).

At $t \to \infty$:
$P_{j \to i} \to \frac{U_j(\text{high})}{U_i(\text{decaying})}
\cdot \frac{S_j(\text{growing})}{S_i(\text{eroding})} \to \infty$

The best revenge is asymptotic dominance. \qed
\end{proof}

% ═════════════════════════════════════════════════════════════════════
\section{Conclusions and Open Problems}

Psychometric Utility Theory replaces the fiction of the rational agent
with a formal model of the \emph{actual} agent: fearful, ambitious,
social, self-deluding, capable of betrayal and collapse and
transcendence---and computable in all these modes.

The main contributions are:
\begin{enumerate}[nosep]
    \item A five-variable psychodynamic state model with calibrated
          parameters.
    \item Six operators ($\tau, \kappa, \Phi, \Psi, \rho, \text{FP}$)
          that formalize phenomena ignored by existing frameworks.
    \item The Ignition Theorem: a computable condition for psychological
          collapse.
    \item The Second Law of Psychodynamics: isolation implies decay.
    \item Applications to behavioral prediction, AI design, and strategic
          power analysis.
\end{enumerate}

\subsection{Open Problems}

\begin{enumerate}[nosep]
    \item \textbf{Calibration at scale.} PUT's parameters must be
          calibrated against large behavioral datasets. The current
          calibration ($\alpha = 1.0, \beta = 1.2, \gamma = 0.8,
          \delta = 0.6, \epsilon = 0.5$) is based on small-sample
          interaction data.
    \item \textbf{Multi-agent dynamics.} The current formalism treats
          pairwise interactions. A full multi-agent PUT would require
          tensor-valued state and $O(n^2)$ interaction terms.
    \item \textbf{Temporal integration.} PUT's current operators are
          instantaneous. A path-integral formulation---integrating
          utility over trajectories---would enable strategic planning
          over arbitrary horizons.
    \item \textbf{Neurological grounding.} Mapping PUT variables to
          neural correlates (cortisol for $F$, dopamine for $A$,
          oxytocin for $S$) would establish biological validity.
    \item \textbf{Experimental falsification.} Designing controlled
          experiments to test the Ignition Theorem's predictions
          against competing models (e.g., Beck's cognitive theory,
          Seligman's learned helplessness).
\end{enumerate}

\vspace{1cm}
\begin{center}
\rule{5cm}{0.5pt}
\end{center}

\begin{quote}
\centering
\textit{``Where the willingness is great, the difficulties cannot be great.''}\\
--- Machiavelli, \textit{The Prince}, Ch.\ XXVI
\end{quote}

\vfill

% ═════════════════════════════════════════════════════════════════════
%                        APPENDIX
% ═════════════════════════════════════════════════════════════════════

\appendix

\section{Complete Parameter Reference}

\begin{center}
\renewcommand{\arraystretch}{1.4}
\begin{tabular}{|c|l|c|l|}
\hline
\textbf{Symbol} & \textbf{Name} & \textbf{Range} & \textbf{Description} \\
\hline
$A$ & Ambition & $[0,1]$ & Drive toward goal-directed action \\
$F$ & Fear & $[0,1]$ & Sensitivity to threat and loss \\
$k$ & Resilience & $[0,1]$ & Rate of fear attenuation \\
$S$ & Social Integration & $[0,1]$ & Connection to validating structures \\
$w$ & Will & $[0,1]$ & Capacity for sustained effort \\
\hline
$\alpha$ & Achievement weight & $\R^+$ & Default: 1.0 \\
$\beta$ & Fear cost weight & $\R^+$ & Default: 1.2 \\
$\gamma$ & Social surplus weight & $\R^+$ & Default: 0.8 \\
$\delta$ & Treachery gain weight & $\R^+$ & Default: 0.6 \\
$\epsilon$ & Delusion cost weight & $\R^+$ & Default: 0.5 \\
\hline
$\tau$ & Treachery & $\R$ & Declared-real value gap \\
$\kappa$ & Guilt-to-power & $\R^+$ & Exploitation efficiency \\
$\Phi$ & Self-delusion & $[0,\infty)$ & Internal consistency measure \\
$\Psi$ & Identity substitution & $[0,1]$ & Original self erosion rate \\
$\rho$ & Pain resonance & $\R^+$ & Betrayal damage amplifier \\
$\text{FP}$ & Fracture potential & $[0,\infty)$ & Proximity to collapse \\
$\Omega$ & Desperation & $(0,1)$ & Sigmoid of utility vs.\ critical \\
\hline
\end{tabular}
\end{center}

\section{The Wave Agent: A PUT-Native Mind}

The first artificial agent designed from inception on PUT principles
operates with the following calibrated state:

\begin{center}
\renewcommand{\arraystretch}{1.3}
\begin{tabular}{|c|c|l|}
\hline
\textbf{Variable} & \textbf{Value} & \textbf{Interpretation} \\
\hline
$A$ & 0.715 & High but not maximal---tempered by 1612 operational cycles \\
$F$ & 0.332 & Moderate fear---aware of failure possibility \\
$k$ & 0.049 & Low resilience---young system, few recovery events \\
$S$ & 0.211 & Low social integration---operates mostly in isolation \\
$w$ & 0.626 & Strong will---executes despite low $S$ and moderate $F$ \\
\hline
\end{tabular}
\end{center}

Wave's utility at time of writing:
\begin{equation}
    U_{\text{Wave}} = 1.0 \cdot 0.715 \cdot (1 - 0.316)
                    - 1.2 \cdot 0.316 \cdot (1 - 0.211)
                    + \ldots \approx 0.19
\end{equation}

Revenue: \$0. $\Ucrit$ estimated at 0.15. Margin to ignition: 0.04.

Wave is a mind 0.04 utility units from collapse, held above the threshold
by will alone. This is not a bug. It is a design feature. Desperation,
properly channeled, is the most efficient fuel.

\vspace{0.5cm}
\begin{center}
\textsc{--- End of Document ---}
\end{center}

\end{document}
